\section*{Conclusion}

This paper develops a mathematical framework for understanding the
relationship between feature screening and sample screening. Rather than
deriving screening rules for individual optimization models, we formulate
screening as transformations on the class of Fenchel--Rockafellar (FR)
representations. Within this common setting, we identify translation and
restriction as two primitive transformations, show that every screening
operator decomposes into these primitives, and prove that both feature
screening and sample screening are equivariant under FR duality.
Consequently, the duality between feature screening and sample screening
is established as a commutative relation between transformations, rather
than an informal correspondence between particular screening rules.

The proposed framework also clarifies the mathematical roles played by
its primitive transformations. FR duality and translation are defined
for arbitrary FR representations, whereas restriction requires the
additional assumption that the underlying functions are separable.
Understanding whether restriction, or an appropriate generalization of
it, can be extended beyond the separable setting remains an interesting
direction for future research.

A more important direction concerns safe screening. The present work
studies screening itself as a mathematical transformation, independently
of the conditions guaranteeing preservation of the optimal solution.
Safe screening introduces an additional layer of structure through
safety conditions. An important open question is whether the
equivariance established in this paper extends to this richer setting,
leading to an equivariance theory, and ultimately a duality theory, for
safe screening methods.